% FINAL-MOVE-5: conclusion — close as screening paper
\section{Conclusion}
\label{sec:conclusion}

We defined frequent losers as a firm-level diagnostic marker and
tested whether it flags procurement environments with higher prices
and greater proximity to known cartels.

It does. The conditional price gap is 3.6--7.7\% across
specifications, and FL-flagged markets overlap with CADE-convicted
cartelists at 3.5$\times$ the baseline rate. Supporting diagnostics
produce results coherent with cover bidding, though none identifies
the mechanism.

We did not identify cover bidding. We did not estimate a causal
effect. Within-market selection remains unresolved: the event-study
pre-trends at $t = -2$ and $t = -3$ preclude a causal reading, and
the FL classification's low temporal persistence (10\% across sample
halves) marks it as an environment indicator, not a firm type. The
Limitations section (Section~\ref{sec:limitations}) catalogs these
and other threats.

The screen is useful precisely because it asks little of the data. It
complements bid-level tools where those tools cannot operate, and it
produces firm-level flags that triage investigations at scale. A flag
is a starting point, not a verdict. Welfare quantification requires a
causal interpretation that this design does not provide; an
illustrative calculation appears in
Appendix~\ref{app:extensions}.

Replication in other procurement systems is the most pressing next
step. The question of whether the price gap is causal remains open
and may require institutional features this setting does not provide.
Whether frequent losers are unlucky competitors or cartel operatives,
their presence marks markets that deserve a closer look.
